% ============================================================
% CHAPTER 3: RELATED WORK
% ============================================================

This chapter locates the thesis in the literature in three passes.
\autoref{sec:related:mad} and \autoref{sec:related:dynamics} survey what
multi-agent debate does and how the interaction dynamics inside it have been
studied; \autoref{sec:related:lenses} weighs the available ways of analysing such
systems against one another; and \autoref{sec:related:gap} states the gap this
thesis addresses.

\subsection{Multi-Agent Debate: Design and Critique}
\label{sec:related:mad}

The multi-agent debate (MAD) paradigm was introduced by
\citet{du2024improving}, who let several instances of a language model
independently propose answers and then critique and revise them over multiple
rounds before committing to a joint response. This ``society of minds'' setup
was reported to improve factual accuracy and mathematical reasoning while
reducing hallucination, treating the underlying model as a black box. The core
intuition, that partially independent agents make partially independent errors
which deliberation can surface and correct, has since motivated a broad family
of systems, from confidence-weighted round tables across heterogeneous models
\citep{chen2024reconcile} to evaluation panels that judge text more reliably
than a single evaluator \citep{chan2024chateval}.

Two design axes recur throughout that family and are the ones we vary.
Communication topology fixes who reads whom: sparse structures can match
fully-connected debate at substantially lower token cost
\citep{li2024sparsetopology}, and across arbitrary networks the placement of a
hub node shifts performance appreciably \citep{liu2025problemsolving}. Memory
fixes how much interaction history each agent retains. The two were recently
crossed in a single factorial by \citet{mehdizadeh2026topologymemory}, who find
no main effect of topology on coordination success but a strong interaction with
memory: longer memory slows settling in decentralised networks and accelerates
it in centralised ones, so the same parameter pushes the system in opposite
directions depending on the network's structure. Faster settling in the centralised
case, they note, means locking into a fragmented plateau rather than reaching
agreement. Beyond these two, the decision protocol itself matters, with voting
favouring reasoning tasks and consensus favouring knowledge tasks, and adding
agents helping where adding rounds can hurt \citep{kaesberg2025voting}.

Against this largely optimistic backdrop, a critical literature has emerged.
\citet{liu2025stopovervaluing} find that MAD frequently fails to outperform
chain-of-thought \citep{wei2022chainofthought} or self-consistency
\citep{wang2023selfconsistency} despite far greater compute, and argue that
model heterogeneity rather than debate is what reliably helps.
\citet{zhang2025debateorvote} show that majority voting alone accounts for most
of the attributed gains, and prove that under their formalisation debate induces
a martingale over belief trajectories and so cannot improve expected correctness
on its own. \citet{estornell2024multillm} supply the mechanism: when agents are
similar in capability or in the responses they produce, debate dynamics become
static and the procedure converges to the majority opinion, so that where the
majority reflects a misconception shared through common training data, debate
converges to the misconception.

Two 2026 results push this further and are the closest published statements of
what we measure. \citet{oh2026entrenchment} decompose the failure into a biased
static initial belief and homogenised dynamics that amplify the majority view
regardless of correctness. \citet{bertalanic2026costofconsensus} show that
unguided homogeneous debate matched or underperformed \emph{isolated
self-correction} at several times the token cost, and that a control injecting
rationales from unrelated problems performed no worse than genuine peer
exchange. They name the resulting loss an \emph{oracle gap}, the distance
between the accuracy achievable from the team's own generation pool and the
accuracy the vote actually returns, a quantity we measure directly in
\autoref{sec:results:perf}.

What unifies the affirmative and the critical strands is that both evaluate
debate almost exclusively by its \emph{final outcome}. The deliberative process
that produces that outcome, and whether two systems reaching the same accuracy do
so through genuinely different dynamics, remains largely outside the scope of
these evaluations.


\subsection{Interaction Dynamics in LLM Collectives}
\label{sec:related:dynamics}

A distinct and more recent line looks inside the debate at how positions evolve
and influence one another. Each of our four motifs corresponds to a debate
phenomenon that this literature documents to a very different degree, and we take
them in turn, before turning to the opening vote distribution, which is not
itself a motif but the input on which they all act.

\paragraph{Rising confidence.}
\citet{prasad2025certain} stage repeated adversarial debates and find pervasive
miscalibration: models begin overconfident, grow \emph{more} confident that they
are winning as the exchange proceeds even absent new supporting evidence, and
frequently both claim victory at once, a logical impossibility. That this drift
is independent of correctness has since been confirmed directly, with
pre-communication confidence a near-worthless predictor of whether an agent is
right, and peer exchange further eroding what little discrimination confidence
carries \citep{bertalanic2026costofconsensus}. Upward, anti-Bayesian and
correctness-blind confidence drift is exactly the kind of positive-feedback
signature that motivates a dynamical treatment.

\paragraph{Conformity and influence concentration.}
\citet{wynn2025talk} show that debate can \emph{decrease} accuracy over time
even when stronger models outnumber weaker ones, as agents abandon correct
answers rather than challenge flawed arguments; they examine sycophancy
\citep{sharma2023sycophancy} and social conformity among the contributing factors. \citet{cho2025herd} find the gap
between an agent's own confidence and its perceived confidence in peers governs
how likely it is to conform, and \citet{mehdizadeh2025peerpressure} characterise
the dose-response as a sigmoid with a model-specific threshold.
\citet{cau2026agreementdrift} identify an aggregate directional asymmetry they
call \emph{agreement drift}, which survives in balanced populations and so cannot
be reduced to a majority effect. Where influence concentrates rather than
diffusing, an emergent leader-follower structure forms
\citep{chen2023consensus}, and outcomes come to depend heavily on the competence
of whichever agent occupies the centre \citep{conformitydynamics2026}, a
dependency our star topology is designed to expose.

\paragraph{Premature lock-in.}
\citet{chuang2025debate} report that role-playing LLM groups converge far more
strongly and rapidly than human groups do, indicating a systematic bias toward
premature agreement. Purely cooperative framings sharpen this: with no dissenting
pressure a group can collapse into a mutually reinforcing echo chamber. A
structurally sharper version is captured by HiddenBench
\citep{li2026hiddenbench}, which adapts the hidden-profile paradigm from social
psychology \citep{stasser1985pooling, schulzhardt2012synergy}: when the evidence
needed for the correct answer is distributed across agents, collective accuracy
collapses far below what a single fully-informed agent achieves, because agents
lock onto shared information and fail to pool their private knowledge. We adopt
HiddenBench as our second benchmark for exactly this reason.

\paragraph{Oscillation and non-termination.}
The fourth regime, sustained oscillation, is conspicuous by its near-absence from
this literature. A debate that never converges but cycles through the same
positions is a limit cycle in the dynamical sense
\citep{strogatz2015nonlinear}, yet debate studies typically impose a fixed round
budget and fold any run that fails to settle into an undifferentiated ``no
consensus'' category, so a periodic orbit is never distinguished from slow
convergence. We flag this near-silence because it prefigures our own finding that
genuine limit cycles are essentially absent outside the star topology
(\autoref{sec:results:lc}).

\paragraph{The opening distribution.}
That round-zero answers govern the outcome is stated in several forms.
\citet{hu2025stabilitydetection} make it a formal premise, proving debate
improvement over majority vote only under the assumption that at least one
initial response is generated under the correct concept.
\citet{zhu2026demystifying} report that initial answer diversity correlates
significantly with final accuracy and propose a diversity-aware initialisation
to raise the chance that a correct hypothesis is present at the start.
Empirically, \citet{wu2025reallydebate} find that an agent starting inside a
correct majority is almost perfectly stable while one facing an incorrect
majority recovers in only a minority of cases (3.6--34.4\% depending on the model). In every case the opening distribution is observed as
it arises rather than constructed, leaving its effect confounded with task
difficulty and agent confidence, which is what \autoref{sec:results:evid} is
designed to separate.

\medskip
\noindent The vocabulary for naming these phenomena precisely, network and
dynamical motifs, positive feedback and lock-in, opinion dynamics and social
power, and the distinction between structural and effective connectivity, was
introduced in \autoref{sec:bg:motifs} and is not repeated here. What that section
leaves open is whether the vocabulary is more than an analogy once it is carried
to LLM populations.

That it is more than an analogy is established empirically.
\citet{ashery2025conventions} demonstrate the spontaneous emergence of
universally adopted conventions in decentralised populations, show that strong
collective biases arise even when individual agents exhibit none, and find that
committed minorities can impose alternative conventions on the majority, the LLM
analogue of the critical-mass threshold measured in human groups
\citep{centola2018tipping}. The same behaviour can be made analytic:
\citet{okawa2026biasedconsensus} derives a mean-field treatment of debate from
spin models, predicting a conformity threshold above which a collective norm
emerges, smoothed into a crossover at the small group sizes typical of debate.


\subsection{Ways of Analysing Multi-Agent Systems}
\label{sec:related:lenses}

Before stating the gap, it is worth making explicit \emph{why} a
complex-systems treatment is the appropriate instrument, by situating it against
the alternatives. Each lens answers a different question and carries different
costs.

\paragraph{Outcome evaluation.}
The dominant approach scores a system by final task performance
\citep{du2024improving, liu2025stopovervaluing, zhang2025debateorvote}. It is
objective, comparable across systems, and measures what practitioners care
about. But an outcome is a scalar summary of a trajectory: two systems reaching
the same accuracy, one by genuine error correction and the other by a confident
slide into a shared mistake, are indistinguishable under it. It says
\emph{whether} a configuration worked, not \emph{why}.

\paragraph{Mechanistic interpretability.}
Opening the models themselves operates at the wrong level of abstraction for our
question. Consensus, oscillation, and influence concentration are properties of
the \emph{group} trajectory rather than of any single forward pass
\citep{miehling2025agentic}, and internals are inaccessible for the API-served
models typical of applied multi-agent systems.

\paragraph{Qualitative transcript analysis.}
Reading transcripts directly is rich and interpretable, and we draw on it
informally when reading the mechanism behind flagged limit-cycle cases
(\autoref{sec:results:lc}). On its own it is subjective, does not scale to tens
of thousands of repetitions, and yields no falsifiable structure testable
against a null.

\paragraph{Influence attribution.}
A fourth approach quantifies directed influence between agents, either
information-theoretically through transfer entropy over the utterances of
interlocutors \citep{infodynamics2026}, or by treating influence as a cooperative
game and estimating it through agent removal \citep{agentsthatmatter2026,
cui2025introspecloo}. These yield causal or information-theoretic estimates
rather than the adoption-based signature we use, at the price of model access or
repeated re-execution; we return to the trade-off in
\autoref{sec:conclusion:future}.

\paragraph{Learned diagnostics from logs.}
A fifth and very recent approach trains a predictor on log-derived features,
either to decide whether a query warrants debate at all \citep{fan2026imad} or
to decide when a majority vote should be overturned
\citep{minoritysentinel2026}. These establish that ordinary logs carry enough
signal to support useful decisions. They are, however, predictive rather than
explanatory: they learn that a feature configuration precedes failure without
naming which dynamical regime the system is in, and they require labelled
outcomes to train.

\paragraph{The dynamical lens.}
The approach taken here treats the debate as a dynamical system and measures
recurring patterns in its state trajectory \citep{strogatz2015nonlinear,
driscoll2024dynamical}. Its advantages match the gaps above: it operates at the
group level, scales to the full dataset, needs only interaction logs and no
labelled outcomes, and gives each motif a falsifiable signature testable against
a principled null. Its disadvantages must be stated as plainly. Transferring
motif concepts from physical and biological systems to an LLM collective is a
modelling choice, not a law; the detection thresholds are partly conventional; a
four-agent group is small for some of these measures; and, absent intervention,
the detectors identify \emph{signatures} and \emph{associations} rather than
certified attractors or causes. We adopt this lens because it is the only one
that can name and measure the pathologies documented above while remaining
applicable to black-box models, but we carry its limitations into the
interpretation of every result (\autoref{sec:discussion:limitations}).


\subsection{Research Gap}
\label{sec:related:gap}

The literature surveyed above has converged on a diagnosis without converging on
an instrument. Outcome-centric evaluation, from the proponents of debate and its
critics alike, measures where a deliberation ends and leaves how it got there
unexamined. The dynamics literature does examine the trajectory, but documents
one pathology at a time, overconfidence here, herding there, premature consensus
elsewhere, and relates them to one another through a shared cause such as
sycophancy rather than through the formal dynamical regimes they resemble. None
of these signatures is tested against a null that says what the same statistic
would look like by chance, so a raw detection rate cannot be distinguished from
the ordinary regularity of a small, discrete state space.

More consequentially, where the complex-systems vocabulary \emph{has} been
applied to LLM collectives, it has been applied on substrates from which truth
was deliberately removed. \citet{ashery2025conventions} study arbitrary
conventions where any label is an acceptable equilibrium.
\citet{mehdizadeh2026topologymemory} verify that their convention symbols carry
no intrinsic model preference, precisely so that dynamics are not confounded by
correctness. \citet{cau2026agreementdrift} choose a non-empirical proposition
explicitly to avoid convergence driven by factual correctness. This is a
principled choice, and it is what makes those studies clean: if the goal is to
understand how a group settles on \emph{a} convention, having no better
convention available is exactly right. It also leaves them silent on the
question a practitioner faces, which is not whether the group locks in but
whether it locks in on the right answer---the epistemic case, where some claims
are better supported than others, remains open.

This thesis takes up that question. The gap it addresses is the absence of an
empirical, literature-grounded detection of self-organization motifs in
\emph{truth-valued} deliberation, and our response follows the three research
questions of \autoref{sec:intro:rqs}. As an \emph{instrument}, we define
detectors for self-reinforcement, multistability, dominance, and limit cycles, each
grounded in the corresponding literature and each equipped with its own null
model, and we run them over a repetition structure that makes
across-repetition attractor statistics computable at all; prior work runs a
single simulation per configuration \citep{cau2026agreementdrift} or analyses
run-level means without across-run endpoint distributions
\citep{mehdizadeh2026topologymemory}, and so cannot measure multistability even
in principle. As a \emph{characterisation}, we apply the suite where every
dynamical quantity can be conditioned on correctness, across two deliberately
contrasting information regimes. As a test of \emph{causation}, we construct the
opening vote composition rather than observing it, separating its effect from
the difficulty and confidence it co-varies with elsewhere. And throughout, we
treat the \emph{absence} of a motif as a finding rather than a failed
measurement, a stance the outcome-centric literature has no natural place for.